% GENERATED FILE. Edit canonical lessons, then run npm run generate:books.
% canonical-source-hash: fnv1a64:accabea5387fe07b
% canonical-lessons: FA-C362-estedlal, FA-C362-estedlal-kardan, FA-C362-manteq, FA-C362-manteqi, FA-C362-qane-kardan

\chapter{Making an argument}
\label{ch:fa-making-an-argument}
\hlchaptermodality{362}

\begin{chapteropening}
\textbf{By the end of this chapter:} \emph{I can make an argument, reason a point through, say whether something is logical and convince someone.}

Complete the last lesson of chapter 362: I can make an argument, reason a point through, say whether something is logical and convince someone.
\end{chapteropening}

\section[\texorpdfstring{\fa{استدلال} (estedlâl)}{استدلال (estedlâl)}]{\fa{استدلال} (estedlâl) — an argument, a line of reasoning}

\label{lesson:FA-C362-estedlal}



\begin{quote}
Before the new one: say the Persian for to believe in, then the Persian for to doubt.
\end{quote}

\subsection*{You'll want to know: \fa{استدلال}}
\textbf{\fa{استدلال}} (\emph{estedlâl}) — \textquotedblleft{}an argument, a line of reasoning\textquotedblright{}.

\begin{quote}
\fa{استدلال} \fa{او} \fa{خیلی} \fa{قوی} \fa{بود}.
\end{quote}

\begin{quote}
\emph{estedlâl-e u kheyli qavi bud.}
\end{quote}

\begin{quote}
His argument was very strong.
\end{quote}

Formal. Plural \emph{estedlâlhâ}. It is an argument as reasoning, not a quarrel.

\begin{cousinweb}
\textbf{\fa{استدلال}} comes from Arabic \emph{istidlāl} \textquotedblleft{}reasoning, inference\textquotedblright{} (EtymDB), a verbal noun from the same Arabic family as \emph{dalīl} \textquotedblleft{}proof, guide\textquotedblright{}, which gave Persian \emph{dalil}.
\end{cousinweb}

\begin{grammarlens}[title={What you've built}]
One more word for telling, explaining and describing at B1.
\end{grammarlens}

\subsection*{Your turn}
\begin{itemize}
\raggedright
  \item \emph{Say it:} \emph{estedlâl}
  \item \emph{Say it:} \emph{estedlâl}, once more
  \item \emph{Say it:} a sentence of your own with \emph{estedlâl}
  \item \emph{Recall:} say the Persian for to believe in, then the Persian for to doubt, then say \emph{estedlâl} again
\end{itemize}

\begin{tcolorbox}[breakable,colback=teal!4,colframe=teal!35!black,title={Before you move on}]
What does \fa{استدلال} mean? (\textquotedblleft{}An argument, a line of reasoning\textquotedblright{}.) Say it once more.
\end{tcolorbox}
\section[\texorpdfstring{\fa{استدلال} \fa{کردن} (estedlâl kardan)}{استدلال کردن (estedlâl kardan)}]{\fa{استدلال} \fa{کردن} (estedlâl kardan) — to argue, to reason}

\label{lesson:FA-C362-estedlal-kardan}



\begin{quote}
Before the new one: say the Persian for to doubt, then the Persian for an argument.
\end{quote}

\subsection*{You'll want to know: \fa{استدلال} \fa{کردن}}
\textbf{\fa{استدلال} \fa{کردن}} (\emph{estedlâl kardan}) — \textquotedblleft{}to argue, to reason\textquotedblright{}.

\begin{quote}
\fa{وکیل} \fa{با} \fa{دقت} \fa{استدلال} \fa{کرد}.
\end{quote}

\begin{quote}
\emph{vakil bâ deqqat estedlâl kard.}
\end{quote}

\begin{quote}
The lawyer argued carefully.
\end{quote}

\emph{Estedlâl} \textquotedblleft{}reasoning\textquotedblright{} + \emph{kardan} (stem \emph{kon-}). To argue that …: \emph{estedlâl mikonad ke} … .

\begin{cousinweb}
\textbf{\fa{استدلال} \fa{کردن}} joins \emph{estedlâl}, from Arabic \emph{istidlāl} \textquotedblleft{}reasoning\textquotedblright{} (EtymDB), to the inherited Persian verb \emph{kardan} \textquotedblleft{}to do\textquotedblright{}.
\end{cousinweb}

\begin{grammarlens}[title={What you've built}]
One more word for telling, explaining and describing at B1.
\end{grammarlens}

\subsection*{Your turn}
\begin{itemize}
\raggedright
  \item \emph{Say it:} \emph{estedlâl kardan}
  \item \emph{Say it:} \emph{estedlâl kardan}, once more
  \item \emph{Say it:} a sentence of your own with \emph{estedlâl kardan}
  \item \emph{Recall:} say the Persian for to doubt, then the Persian for an argument, then say \emph{estedlâl kardan} again
\end{itemize}

\begin{tcolorbox}[breakable,colback=teal!4,colframe=teal!35!black,title={Before you move on}]
What does \fa{استدلال} \fa{کردن} mean? (\textquotedblleft{}To argue, to reason\textquotedblright{}.) Say it once more.
\end{tcolorbox}
\section[\texorpdfstring{\fa{منطق} (manteq)}{منطق (manteq)}]{\fa{منطق} (manteq) — logic}

\label{lesson:FA-C362-manteq}



\begin{quote}
Before the new one: say the Persian for an argument, then the Persian for to argue.
\end{quote}

\subsection*{You'll want to know: \fa{منطق}}
\textbf{\fa{منطق}} (\emph{manteq}) — \textquotedblleft{}logic\textquotedblright{}.

\begin{quote}
\fa{حرف} \fa{تو} \fa{منطق} \fa{ندارد}.
\end{quote}

\begin{quote}
\emph{harf-e to manteq nadârad.}
\end{quote}

\begin{quote}
What you say has no logic.
\end{quote}

Also the subject of logic at university.

\begin{cousinweb}
\textbf{\fa{منطق}} was borrowed from Arabic \emph{manṭiq} (EtymDB glosses it \textquotedblleft{}speech\textquotedblright{}); in Arabic and Persian it came to mean logic, the art of correct speech and reasoning.
\end{cousinweb}

\begin{grammarlens}[title={What you've built}]
One more word for telling, explaining and describing at B1.
\end{grammarlens}

\subsection*{Your turn}
\begin{itemize}
\raggedright
  \item \emph{Say it:} \emph{manteq}
  \item \emph{Say it:} \emph{manteq}, once more
  \item \emph{Say it:} a sentence of your own with \emph{manteq}
  \item \emph{Recall:} say the Persian for an argument, then the Persian for to argue, then say \emph{manteq} again
\end{itemize}

\begin{tcolorbox}[breakable,colback=teal!4,colframe=teal!35!black,title={Before you move on}]
What does \fa{منطق} mean? (\textquotedblleft{}Logic\textquotedblright{}.) Say it once more.
\end{tcolorbox}
\section[\texorpdfstring{\fa{منطقی} (manteqi)}{منطقی (manteqi)}]{\fa{منطقی} (manteqi) — logical, reasonable}

\label{lesson:FA-C362-manteqi}



\begin{quote}
Before the new one: say the Persian for to argue, then the Persian for logic.
\end{quote}

\subsection*{You'll want to know: \fa{منطقی}}
\textbf{\fa{منطقی}} (\emph{manteqi}) — \textquotedblleft{}logical, reasonable\textquotedblright{}.

\begin{quote}
\fa{این} \fa{یک} \fa{تصمیم} \fa{منطقی} \fa{است}.
\end{quote}

\begin{quote}
\emph{in yek tasmim-e manteqi ast.}
\end{quote}

\begin{quote}
This is a logical decision.
\end{quote}

\emph{Manteq} + the adjective suffix \emph{-i}. The opposite is \emph{gheyr-e manteqi}.

\begin{cousinweb}
\textbf{\fa{منطقی}} comes from Arabic \emph{manṭiqī} \textquotedblleft{}logical\textquotedblright{} (EtymDB), the adjective of \emph{manṭiq} \textquotedblleft{}logic\textquotedblright{}; Persian hears it as \emph{manteq} + \emph{-i}.
\end{cousinweb}

\begin{grammarlens}[title={What you've built}]
One more word for telling, explaining and describing at B1.
\end{grammarlens}

\subsection*{Your turn}
\begin{itemize}
\raggedright
  \item \emph{Say it:} \emph{manteqi}
  \item \emph{Say it:} \emph{manteqi}, once more
  \item \emph{Say it:} a sentence of your own with \emph{manteqi}
  \item \emph{Recall:} say the Persian for to argue, then the Persian for logic, then say \emph{manteqi} again
\end{itemize}

\begin{tcolorbox}[breakable,colback=teal!4,colframe=teal!35!black,title={Before you move on}]
What does \fa{منطقی} mean? (\textquotedblleft{}Logical, reasonable\textquotedblright{}.) Say it once more.
\end{tcolorbox}
\section[\texorpdfstring{\fa{قانع} \fa{کردن} (qâne' kardan)}{قانع کردن (qâne' kardan)}]{\fa{قانع} \fa{کردن} (qâne' kardan) — to convince, to persuade}

\label{lesson:FA-C362-qane-kardan}



\begin{quote}
Before the new one: say the Persian for logic, then the Persian for logical.
\end{quote}

\subsection*{You'll want to know: \fa{قانع} \fa{کردن}}
\textbf{\fa{قانع} \fa{کردن}} (\emph{qâne' kardan}) — \textquotedblleft{}to convince, to persuade\textquotedblright{}.

\begin{quote}
\fa{بالاخره} \fa{دوستم} \fa{را} \fa{قانع} \fa{کردم}.
\end{quote}

\begin{quote}
\emph{bâlâkhare dustam râ qâne' kardam.}
\end{quote}

\begin{quote}
In the end I convinced my friend.
\end{quote}

\emph{Qâne'} \textquotedblleft{}satisfied, content\textquotedblright{} + \emph{kardan}. \emph{Qâne' shodan} — to be convinced. The person takes \emph{râ}.

\begin{cousinweb}
\textbf{\fa{قانع}} is usually traced to Arabic \emph{qāni'} \textquotedblleft{}content, satisfied\textquotedblright{}; with the Persian verb \emph{kardan} \textquotedblleft{}to do\textquotedblright{}, to convince someone is to make them content with your view.
\end{cousinweb}

\begin{grammarlens}[title={What you've built}]
One more word for telling, explaining and describing at B1.
\end{grammarlens}

\subsection*{Your turn}
\begin{itemize}
\raggedright
  \item \emph{Say it:} \emph{qâne' kardan}
  \item \emph{Say it:} \emph{qâne' kardan}, once more
  \item \emph{Say it:} a sentence of your own with \emph{qâne' kardan}
  \item \emph{Recall:} say the Persian for logic, then the Persian for logical, then say \emph{qâne' kardan} again
\end{itemize}

\begin{tcolorbox}[breakable,colback=teal!4,colframe=teal!35!black,title={Before you move on}]
What does \fa{قانع} \fa{کردن} mean? (\textquotedblleft{}To convince, to persuade\textquotedblright{}.) Say it once more.
\end{tcolorbox}
